% Generated by scripts/assemble.py; corrections belong in source/corrections.json.
\examhead{121}{山东大学 2024 年数学分析试卷}{116}
% p117-b10; source PDF page 117
\par\Needspace{5\baselineskip}
\blockanchor{p117-b10}
\majorquestion{一}计算题. 每题 15 分, 共 60 分.\par
% p117-b11; source PDF page 117
\blockanchor{p117-b11}
\question{1}求 $\sum_{n=1}^{\infty} \frac{x + n(-1)^n}{x^2 + n^2}$ 的条件收敛域, 绝对收敛域, 以及一致收敛域.\par
% p117-b12; source PDF page 117
\blockanchor{p117-b12}
\question{2}比较下列无穷大量.\par
% p117-b13; source PDF page 117
\blockanchor{p117-b13}
(1) $x$ 与 $(\ln x)^{100}(x \to +\infty)$.\par
% p117-b14; source PDF page 117
\blockanchor{p117-b14}
(2) $(\ln x)^{100}$ 与 $e^{(\ln x)^{\frac{1}{100}}} (x \to +\infty)$.\par
% p117-b15; source PDF page 117
\blockanchor{p117-b15}
\question{3}计算第一类曲面积分 $\iint_{\Sigma} (x + y + z) \, \mathrm{d}S$, 其中曲面 $\Sigma$ 是左半球面 $x^2 + y^2 + z^2 = 1, y \leq 0$.\par
% p117-b16; source PDF page 117
\blockanchor{p117-b16}
\question{4}计算 $I(\alpha) = \int_{0}^{+\infty} \frac{\arctan(\alpha x)}{x(1 + x^2)} \, \mathrm{d}x$ ($\alpha > 0$).\par
% p117-b17; source PDF page 117
\par\Needspace{5\baselineskip}
\blockanchor{p117-b17}
\majorquestion{二}证明题. 每题 20 分, 共 60 分.\par
% p117-b18; source PDF page 117
\blockanchor{p117-b18}
\question{1}设  $f_{n}(x)=1+\frac{1}{2^{x}}+\frac{1}{3^{x}}+\cdots+\frac{1}{n^{x}}$ ，证明：对任何  $\alpha>1$ ，函数列  $\{f_{n}(x)\}$  在  $[\alpha,+\infty)$  上一致收敛.\par
% p117-b19; source PDF page 117
\blockanchor{p117-b19}
\question{2}解答如下问题:\par
% p117-b20; source PDF page 117
\blockanchor{p117-b20}
(1) 证明不等式 $\frac{x}{1 + x} < \ln (1 + x) < x (x > 0)$.\par
% p117-b21; source PDF page 117
\blockanchor{p117-b21}
(2) 利用上述不等式证明数列  $\{a_{n}\}$  存在极限, 其中  $a_{n}=1+\frac{1}{2}+\frac{1}{3}+\cdots+\frac{1}{n}-\ln n$ .\par
% p117-b22; source PDF page 117
\par\Needspace{6\baselineskip}
\blockanchor{p117-b22}
\question{3}设 $f(x)$ 在区间 $[0,1]$ 上连续, 证明:\par
\nopagebreak[4]
% p117-b23; source PDF page 117
\blockanchor{p117-b23}
\displayformula{\int_ {0} ^ {\pi} x f (\sin x) \mathrm{d} x = \frac {\pi}{2} \int_ {0} ^ {\pi} f (\sin x) \mathrm{d} x.}
% p117-b24; source PDF page 117
\blockanchor{p117-b24}
并由此计算 $I = \int_0^\pi x\frac{\sin^{2n}x}{\sin^{2n}x + \cos^{2n}x}\mathrm{d}x.$\par
% p117-b25; source PDF page 117
\par\Needspace{5\baselineskip}
\blockanchor{p117-b25}
\majorquestion{三}综合题. 每题 15 分, 共 30 分.\par
% p117-b26; source PDF page 117
\blockanchor{p117-b26}
\question{1}设 $z = z(x, y)$ 满足 $F\left(x + \frac{z}{y}, y + \frac{z}{x}\right) = 0$，其中 $z(x, y), F(u, v)$ 均可微，求 $z - xz_x' - yz_y'$.\par
% p117-b27; source PDF page 117
\blockanchor{p117-b27}
\question{2}设二元连续可微函数 $F$ 在直角坐标下可写为 $F(x, y) = f(x)g(y)$, 在极坐标下可写为 $F(r \cos \theta, r \sin \theta) = h(r)$, 若 $F(x, y)$ 无零点, 求 $F(x, y)$.\par
